% YAML frontmatter for the website sits in the block below. LaTeX ignores it
% (it is a comment); the converter lifts it verbatim into the module
% frontmatter and CI validates it (pure YAML).
% Keys: title, summary, contributors — all optional. The website extracts
% title/contributors from \title{} and \author{}; a frontmatter key here
% overrides the extracted value. `summary:` is the sheet's OVERVIEW — the
% page header shows it under the title, and it doubles as the index blurb,
% so it lives here (metadata, never an "Overview" body section) and stays
% one tight paragraph. Use a YAML folded block scalar (`>-`) with
% continuation lines indented two spaces after the leading `% `.
% Where the sheet sits in the course — its cluster and teaching day — is NOT a
% key here: list the slug under its day in schedule.yaml and the build stamps
% both in. (`unlisted: true` is a maintainer flag: the page is built but
% linked from nowhere, and excused from the schedule — this demo is the one
% sheet that wants it.)
%--- iliad -------------------------------------------------------------------
% title: "Example Worksheet: The Kitchen Sink"
% summary: >-
%   A dummy worksheet demonstrating every construct the pipeline supports.
% unlisted: true
%--- end ----------------------------------------------------------------------
\documentclass[11pt]{article}

% The ONE required package. It provides: exercise, solution, proof, callout,
% remark, definition/theorem/lemma/proposition/corollary/fact/example,
% the \ifsolutions toggle, and loads hyperref+cleveref (configure hyperref
% with \hypersetup{...}; don't re-load it). PDF styling is plain by default —
% style your PDF however you like, or pass [boxes] for ready-made colour boxes.
% iliad.sty: local copy first (Overleaf), else the shared ../iliad.sty
\IfFileExists{iliad.sty}{\usepackage[boxes]{iliad}}{\usepackage[boxes]{../iliad}}

% \solutionsfalse   % uncomment to hide all solutions from the PDF

% ---- Free zone -------------------------------------------------------------
% Define whatever notation you like; load extra packages as needed.
% Macros become KaTeX \gdef's automatically (avoid \mathchoice / optional-arg
% macros — the converter warns if it can't translate one).
\usepackage{tikz}
\newcommand{\R}{\mathbb{R}}
\newcommand{\E}{\mathbb{E}}
\newcommand{\KL}[2]{D_{\mathrm{KL}}\!\left(#1 \,\middle\|\, #2\right)}
\DeclareMathOperator*{\argmax}{arg\,max}

\title{Example Worksheet: The Kitchen Sink}
\author{\authorname{Jane Doe}\\ \affiliation{Example University}
  \and John Roe}
\date{\today}

\begin{document}
\maketitle

% Front matter opens every sheet, in this fixed order (docs/commands.md
% §"Front matter opens the sheet"):
%   1. video embeds (optional)   2. Prerequisites   3. the learning-outcomes
% box — then the content. The sheet's OVERVIEW is not a body section: it is
% the `summary:` in the metadata block up top, and the page header shows it
% under the title. The build prints a non-fatal warning when a sheet strays
% from this order or keeps an Overview section in the body.

% \youtube[Title]{VIDEO_ID} — pass the 11-character video ID (the v= value of
% the watch URL), never the full URL. Embedded player on the web; here in the
% PDF, a "Video:" line carrying the watch URL. With no [Title], the web build
% queries the title from YouTube itself. Optional — most sheets have none.
% Lecture recordings go here, first; a video illustrating one point may sit
% with that point instead.
\youtube{dQw4w9WgXcQ}

%===============================================================================
\section{Prerequisites}
\label{sec:prereqs}

Prerequisites, further reading, and anything else you want to say are just
ordinary \LaTeX{} — no special markup. For this sheet:
\begin{itemize}
  \item comfort with probability on finite sets ($\E$, distributions);
  \item one variable of calculus (we differentiate $-t \log t$ once).
\end{itemize}

% The "What you'll learn" box: after Prerequisites, before the content.
% The body is ordinary LaTeX. Simplest form: a single itemize.
%   \begin{learningoutcomes}
%     \begin{itemize}
%       \item ...
%     \end{itemize}
%   \end{learningoutcomes}
% For a longer sheet, group the outcomes under \subsection*{...} headings,
% each followed by its own itemize (the headings become bold subheadings in
% the box, on both the PDF and the web) -- shown here:
\begin{learningoutcomes}
  \subsection*{Using this template}
  \begin{itemize}
    \item Recognise every markup element the tex $\to$ MDX pipeline supports.
    \item Copy this file as the starting point for a new worksheet.
  \end{itemize}

  \subsection*{Structuring learning outcomes}
  \begin{itemize}
    \item Group outcomes with \verb|\subsection*{...}| plus an \verb|itemize| when a flat list gets long.
    \item Drop the \verb|\subsection*| headings for a short sheet -- a single \verb|itemize| is fine.
  \end{itemize}
\end{learningoutcomes}

This dummy worksheet demonstrates every construct that survives the round trip
from \LaTeX{} to the curriculum website. Compile it with \texttt{pdflatex}
(twice, then \texttt{bibtex}, then twice again for the bibliography), or just
read the source side by side with the rendered page.

%===============================================================================
\section{Prose, math, and cross-references}
\label{sec:basics}

Ordinary prose converts verbatim, including \emph{emphasis}, \textbf{bold},
\texttt{monospace}, and special characters like 50\% or \S3. Inline math is
copied byte-for-byte: $e^{i\pi} + 1 = 0$, and KaTeX handles comparisons like
$x_{<t}$ without trouble. Custom notation defined in the preamble just works:
$\KL{p}{q} \geq 0$, $\argmax_{a} \E[R \mid a]$, for $p, q$ distributions on
$\R^n$.

Display math works as \, \[ \sum_{k=1}^{n} k = \frac{n(n+1)}{2}, \]
and as numbered equations you can reference:
\begin{equation}
  \label{eq:basel}
  \sum_{n \geq 1} \frac{1}{n^2} = \frac{\pi^2}{6}.
\end{equation}
Multi-line derivations use \texttt{align}:
\begin{align}
  (a+b)^2 &= (a+b)(a+b) \\
          &= a^2 + 2ab + b^2.
\end{align}
Cross-references resolve to the exact text LaTeX prints: \cref{eq:basel} sits
in \cref{sec:basics}; below we will meet \cref{def:entropy,thm:gibbs}. Lists
are fine too:
\begin{itemize}
  \item first bullet;
  \item second bullet, with math $2^n$.
\end{itemize}

%===============================================================================
\section{Definitions, theorems, and callouts}
\label{sec:results}

Label anything you want to reference. Any placement \LaTeX{} binds correctly
is valid (top level of the environment); right after the \texttt{begin} is
clearest.

\begin{definition}[entropy]
\label{def:entropy}
The \emph{entropy} of a distribution $p$ on a finite set $\mathcal{X}$ is
\[ H(p) = -\sum_{x \in \mathcal{X}} p(x) \log p(x). \]
\end{definition}

\begin{theorem}[Gibbs' inequality]
\label{thm:gibbs}
For any distributions $p, q$ on $\mathcal{X}$,
$-\sum_x p(x) \log q(x) \geq H(p)$, with equality iff $p = q$.
\end{theorem}

\begin{proof}
By definition, $\KL{p}{q} = \sum_x p(x) \log \frac{p(x)}{q(x)} \geq 0$ by
Jensen's inequality applied to the concave function $\log$; rearranging gives
the claim. Equality in Jensen requires $p(x)/q(x)$ constant, i.e.\ $p = q$.
(On the website this proof is collapsible, like a solution.)
\end{proof}

\begin{lemma}
\label{lem:uniform}
$H(p) \leq \log |\mathcal{X}|$, with equality iff $p$ is uniform.
\end{lemma}

\begin{proof}
Apply \cref{thm:gibbs} with $q$ uniform.
\end{proof}

\begin{callout}[warning]
\label{co:pitfall}
Callouts flag pitfalls, tips, or side notes. Types: \texttt{note},
\texttt{tip}, \texttt{warning}. They may be labelled and referenced,
like this one.
\end{callout}

\begin{remark}
\label{rem:aside}
A \texttt{remark} is an aside; on the website it renders as a note-style
callout. Referencing works: see \cref{co:pitfall} above (and this remark is
\cref{rem:aside}).
\end{remark}

%===============================================================================
\section{Exercises and solutions}
\label{sec:exercises}

The optional argument of an exercise is its title, exactly as in amsthm.
An optional \verb|\important| mark (a $\bigstar$ flagging the sheet's most
important exercises) goes right after the label. Subparts are a plain
\texttt{enumerate}; label an item to reference it. Every solution
\emph{names its exercise} — \texttt{[ex:label]} is mandatory — so you may
place solutions right after each exercise or collect them at the end, as you
prefer. Exercises don't require solutions; label any exercise a solution or
reference points at.

\begin{exercise}[Entropy warm-up]
\label{ex:warmup}
Let $p$ be a distribution on a finite set $\mathcal{X}$.
\begin{enumerate}
  \item Show that $H(p) \geq 0$.
        \label{ex:warmup-a}
  \item For which $p$ is $H(p) = 0$?
\end{enumerate}
\end{exercise}

\begin{solution}[ex:warmup]
For \cref{ex:warmup-a}: each term $-p(x)\log p(x) \geq 0$ since
$p(x) \in [0,1]$. For part (b): $H(p) = 0$ iff every term vanishes iff $p$ is
a point mass. \hint{when is $-t \log t = 0$ on $[0,1]$?}
\end{solution}

\begin{exercise}[Gibbs, the hard direction]
\label{ex:gibbs-hard}
\important
Prove the equality case of \cref{thm:gibbs} directly, without citing Jensen's
equality condition. \note{this one is solved at the end of the sheet instead
of inline, to demonstrate both solution placements.}
\end{exercise}

Hints also have a block form: unnumbered, never labeled, rendered right where
it is written — a bold lead-in on paper, a collapsible drop-down on the web.

\begin{hint}
Consider $f(t) = t - 1 - \log t$: it is nonnegative on $t > 0$ and vanishes
only at $t = 1$.
\end{hint}

%===============================================================================
\section{Figures and citations}
\label{sec:figures}

Figures use \texttt{figure} + \verb|\includegraphics| (export TikZ to PDF the
same way; keep assets in \texttt{fig/}).

\begin{figure}[ht]
  \centering
  \includegraphics[width=0.45\linewidth]{fig/value-curve.pdf}
  \caption{An example figure. On the web this becomes a \texttt{Figure}
  component.}
  \label{fig:value}
\end{figure}

See \cref{fig:value}. Citations come from the \texttt{.bib} file: the idea of
entropy is due to \cite{Shannon:48}; for a textbook treatment see
\cite[ch.~4]{MacKay:03}. External links are plain
\href{https://en.wikipedia.org/wiki/Entropy_(information_theory)}{hyperlinks},
and footnotes are supported.\footnote{On the web this becomes a numbered
  footnote: the marker links down to a note at the foot of the page, and the
  note links back. In the PDF it is an ordinary \LaTeX{} footnote.}

%===============================================================================
\section*{Selected solutions}

\begin{solution}[ex:gibbs-hard]
Write $f(t) = t - 1 - \log t$ for $t > 0$ and note $f(t) \geq 0$ with equality
iff $t = 1$. Then with $t_x = q(x)/p(x)$,
\[ -\sum_x p(x) \log q(x) - H(p) = \sum_x p(x) f\!\left(t_x\right) +
   \sum_x \bigl(p(x) - q(x)\bigr) \geq 0, \]
and equality forces $t_x = 1$ wherever $p(x) > 0$, i.e.\ $p = q$.
\end{solution}

% A solutionsonly block appears ONLY in the with-solutions build: it renders as
% plain content here and in the with-solutions PDF, and is removed entirely from
% every -nosol download (like a solution, but with no box, heading, or exercise
% binding). Use it for an answer key, an instructor aside, or any spoiler.
\begin{solutionsonly}
\textbf{Instructor note.} \cref{ex:gibbs-hard} is the one to discuss on
the board if time is short; the warm-up can be left for self-study.
\end{solutionsonly}

\bibliographystyle{plain}
\bibliography{biblo}

\end{document}
